
\subsubsection{2.1 Spurious Correlation Detection and Mitigation}

Sagawa et al. introduced Group DRO, which minimizes worst-group loss when group annotations are available, establishing the Waterbirds benchmark. Annotation-free methods have emerged: JTT exploits the observation that models misclassify minority-group examples early in training; Environment Inference for Invariant Learning clusters training examples to infer pseudo-environments; Spread Spurious Attribute estimates spurious attributes without supervision.

These methods implicitly rely on early-training signals to identify spurious reliance but do not directly measure the gradient dynamics underlying why spurious features are learned preferentially. Our work tests the mechanistic assumption that spurious features receive stronger gradient signals.

\subsubsection{2.2 Simplicity Bias and Learning Dynamics}

Shah et al. formalized simplicity bias: neural networks preferentially learn simpler features when multiple predictive features are available. Their theoretical analysis shows that linear networks and nonlinear networks with certain activation functions exhibit this bias due to gradient flow properties. However, they study feature learning order rather than per-sample gradient magnitudes.

Dataset Cartography maps training dynamics by tracking prediction confidence and variability across epochs, identifying ''easy,'' ''hard,'' and ''ambiguous'' examples. While this characterizes learning dynamics, it does not directly connect to spurious feature attribution or test the gradient competition hypothesis.

Arpit et al. showed that networks learn ''easy'' patterns before memorizing noise, and Kalimeris et al. demonstrated increasing complexity of learned functions over training. These works establish that learning order depends on pattern simplicity but do not measure whether simpler patterns receive stronger or weaker gradient signals.

Our work bridges this gap by directly measuring gradient norms for spurious-aligned versus minority samples.

\subsubsection{2.3 Optimization-Based Robustness}

Sharpness-Aware Minimization (SAM) seeks parameters in flat loss landscape regions, improving generalization across domains. The connection between loss landscape geometry and spurious feature reliance remains underexplored. Our falsification of the gradient competition hypothesis suggests that convergence speed---how quickly different samples reach low loss---may be more relevant than gradient magnitude. This connects to SAM's potential mechanism: if spurious patterns occupy sharper minima, SAM's sharpness penalty would disfavor them.
